\documentclass[12pt]{article}

% Packages
\usepackage[
  paperwidth=8.5in,
  paperheight=11in,
  margin=0.5in
]{geometry}
\usepackage{amsmath, amssymb}
\usepackage{booktabs}
\usepackage{tabularx}
\usepackage{graphicx}
\usepackage{hyperref}


\newcommand{\tacred}{FS-TACRED}
\newcommand{\qwensmall}{Qwen3-4B}
\newcommand{\gemmasmall}{Gemma3-4B}

\begin{document}

Here is the short overview of the experiment that I explored this week:\\

The goal is to use both "natural language based gradient descent" and "pure gradients based" prompt  optimization in the pipeline.
The idea is to optimize the prompt using the former one at first since it can optimize quickly (i.e. using few iterations it can reach near optimal region), 
then apply pure gradient based optimization (smaller and slower improvements but more likely to improve at the near optimal region).

\paragraph{Gradient based optimization.} We first locate some regions (a series of tokens, usually length of 1) where the gradients are highest.
Then generate some candidate replacements of those regions by two ways: 1) prompting LLM to suggest simplest replacements 2) use LM probability to generate candidates for those region.
After we get all the candidates, we place it into the prompt and validate the prompt.

I am still running the experiments. Hopefully, in the meeting, I can show the preliminary results and discuss the method in more details.


% \bibliographystyle{plain}  
% \bibliography{custom}    

\end{document}
